\documentclass[11pt,a4paper]{article}
\newcommand{\intestazione}{Trimero ad anello --- Parte 2, il filo conduttore}
% =====================================================================
%  Preambolo comune ai documenti sul dimero (serie "che cosa abbiamo fatto")
%  Incluso con % =====================================================================
%  Preambolo comune ai documenti sul dimero (serie "che cosa abbiamo fatto")
%  Incluso con % =====================================================================
%  Preambolo comune ai documenti sul dimero (serie "che cosa abbiamo fatto")
%  Incluso con \input{preambolo_dimero} subito dopo \documentclass.
% =====================================================================
\usepackage[T1]{fontenc}
\usepackage[utf8]{inputenc}
\usepackage[main=italian,provide=*]{babel}
\usepackage{amsmath,amssymb,mathtools}
\usepackage{graphicx}
\usepackage{booktabs}
\usepackage{colortbl}
\usepackage{xcolor}
\usepackage{geometry}
\usepackage{placeins}
\usepackage{caption}
\usepackage{enumitem}
\usepackage{fancyhdr}
\usepackage{needspace}
\usepackage{tikz}
\usetikzlibrary{arrows.meta,positioning,shapes.geometric,calc,fit,backgrounds}
\usepackage{hyperref}
\hypersetup{colorlinks=true,linkcolor=blu,citecolor=blu,urlcolor=blu}

\geometry{margin=2.5cm}

% --- colori coerenti con quelli delle figure (palette Okabe-Ito)
\definecolor{blu}{HTML}{0072B2}
\definecolor{arancio}{HTML}{E69F00}
\definecolor{verdes}{HTML}{009E73}
\definecolor{rossos}{HTML}{D55E00}
\definecolor{violas}{HTML}{CC79A7}
\definecolor{grigetto}{HTML}{E8E8E8}

\captionsetup{font=small,labelfont=bf,width=0.92\textwidth}

\pagestyle{fancy}
\fancyhf{}
\fancyhead[L]{\small\itshape\intestazione}
\fancyhead[R]{\small\thepage}
\renewcommand{\headrulewidth}{0.3pt}

% --- notazione
\newcommand{\ket}[1]{\lvert #1\rangle}
\newcommand{\bra}[1]{\langle #1 \rvert}
\newcommand{\media}[1]{\langle #1 \rangle}
\newcommand{\vs}[1]{\mathbf{s}_{#1}}

% --- riquadro "in una riga": il messaggio da portare a casa
\newcommand{\messaggio}[1]{%
  \begin{center}
  \begin{minipage}{0.93\textwidth}
  \setlength{\fboxsep}{9pt}%
  \colorbox{grigetto}{\begin{minipage}{\dimexpr\textwidth-2\fboxsep}
  \small\textbf{In una riga.}\enspace #1
  \end{minipage}}
  \end{minipage}
  \end{center}
}

% --- riquadro "come lo sappiamo": la verifica numerica dietro un'affermazione
\newcommand{\verifica}[1]{%
  \begin{center}
  \begin{minipage}{0.93\textwidth}
  \small\textcolor{verdes}{\rule{2.2em}{1.1pt}}\enspace
  \textcolor{verdes}{\textbf{Come lo sappiamo.}}\enspace #1
  \end{minipage}
  \end{center}
  \medskip
}
 subito dopo \documentclass.
% =====================================================================
\usepackage[T1]{fontenc}
\usepackage[utf8]{inputenc}
\usepackage[main=italian,provide=*]{babel}
\usepackage{amsmath,amssymb,mathtools}
\usepackage{graphicx}
\usepackage{booktabs}
\usepackage{colortbl}
\usepackage{xcolor}
\usepackage{geometry}
\usepackage{placeins}
\usepackage{caption}
\usepackage{enumitem}
\usepackage{fancyhdr}
\usepackage{needspace}
\usepackage{tikz}
\usetikzlibrary{arrows.meta,positioning,shapes.geometric,calc,fit,backgrounds}
\usepackage{hyperref}
\hypersetup{colorlinks=true,linkcolor=blu,citecolor=blu,urlcolor=blu}

\geometry{margin=2.5cm}

% --- colori coerenti con quelli delle figure (palette Okabe-Ito)
\definecolor{blu}{HTML}{0072B2}
\definecolor{arancio}{HTML}{E69F00}
\definecolor{verdes}{HTML}{009E73}
\definecolor{rossos}{HTML}{D55E00}
\definecolor{violas}{HTML}{CC79A7}
\definecolor{grigetto}{HTML}{E8E8E8}

\captionsetup{font=small,labelfont=bf,width=0.92\textwidth}

\pagestyle{fancy}
\fancyhf{}
\fancyhead[L]{\small\itshape\intestazione}
\fancyhead[R]{\small\thepage}
\renewcommand{\headrulewidth}{0.3pt}

% --- notazione
\newcommand{\ket}[1]{\lvert #1\rangle}
\newcommand{\bra}[1]{\langle #1 \rvert}
\newcommand{\media}[1]{\langle #1 \rangle}
\newcommand{\vs}[1]{\mathbf{s}_{#1}}

% --- riquadro "in una riga": il messaggio da portare a casa
\newcommand{\messaggio}[1]{%
  \begin{center}
  \begin{minipage}{0.93\textwidth}
  \setlength{\fboxsep}{9pt}%
  \colorbox{grigetto}{\begin{minipage}{\dimexpr\textwidth-2\fboxsep}
  \small\textbf{In una riga.}\enspace #1
  \end{minipage}}
  \end{minipage}
  \end{center}
}

% --- riquadro "come lo sappiamo": la verifica numerica dietro un'affermazione
\newcommand{\verifica}[1]{%
  \begin{center}
  \begin{minipage}{0.93\textwidth}
  \small\textcolor{verdes}{\rule{2.2em}{1.1pt}}\enspace
  \textcolor{verdes}{\textbf{Come lo sappiamo.}}\enspace #1
  \end{minipage}
  \end{center}
  \medskip
}
 subito dopo \documentclass.
% =====================================================================
\usepackage[T1]{fontenc}
\usepackage[utf8]{inputenc}
\usepackage[main=italian,provide=*]{babel}
\usepackage{amsmath,amssymb,mathtools}
\usepackage{graphicx}
\usepackage{booktabs}
\usepackage{colortbl}
\usepackage{xcolor}
\usepackage{geometry}
\usepackage{placeins}
\usepackage{caption}
\usepackage{enumitem}
\usepackage{fancyhdr}
\usepackage{needspace}
\usepackage{tikz}
\usetikzlibrary{arrows.meta,positioning,shapes.geometric,calc,fit,backgrounds}
\usepackage{hyperref}
\hypersetup{colorlinks=true,linkcolor=blu,citecolor=blu,urlcolor=blu}

\geometry{margin=2.5cm}

% --- colori coerenti con quelli delle figure (palette Okabe-Ito)
\definecolor{blu}{HTML}{0072B2}
\definecolor{arancio}{HTML}{E69F00}
\definecolor{verdes}{HTML}{009E73}
\definecolor{rossos}{HTML}{D55E00}
\definecolor{violas}{HTML}{CC79A7}
\definecolor{grigetto}{HTML}{E8E8E8}

\captionsetup{font=small,labelfont=bf,width=0.92\textwidth}

\pagestyle{fancy}
\fancyhf{}
\fancyhead[L]{\small\itshape\intestazione}
\fancyhead[R]{\small\thepage}
\renewcommand{\headrulewidth}{0.3pt}

% --- notazione
\newcommand{\ket}[1]{\lvert #1\rangle}
\newcommand{\bra}[1]{\langle #1 \rvert}
\newcommand{\media}[1]{\langle #1 \rangle}
\newcommand{\vs}[1]{\mathbf{s}_{#1}}

% --- riquadro "in una riga": il messaggio da portare a casa
\newcommand{\messaggio}[1]{%
  \begin{center}
  \begin{minipage}{0.93\textwidth}
  \setlength{\fboxsep}{9pt}%
  \colorbox{grigetto}{\begin{minipage}{\dimexpr\textwidth-2\fboxsep}
  \small\textbf{In una riga.}\enspace #1
  \end{minipage}}
  \end{minipage}
  \end{center}
}

% --- riquadro "come lo sappiamo": la verifica numerica dietro un'affermazione
\newcommand{\verifica}[1]{%
  \begin{center}
  \begin{minipage}{0.93\textwidth}
  \small\textcolor{verdes}{\rule{2.2em}{1.1pt}}\enspace
  \textcolor{verdes}{\textbf{Come lo sappiamo.}}\enspace #1
  \end{minipage}
  \end{center}
  \medskip
}


\title{\bfseries Il trimero ad anello come sistema aperto\\[2pt]
\large Documento 0 --- il filo conduttore della Parte 2}
\author{Samuele Schiavi \\ \small Tesi triennale in Fisica, Università di Parma
--- Relatore: Prof.\ Alessandro Chiesa}
\date{}

\begin{document}
\maketitle
\thispagestyle{fancy}

\begin{abstract}
\noindent
Questa serie racconta l'estensione al trimero ad anello della Parte 2 già
completata sul dimero: introdurre il rumore, passare dal vettore di stato
all'operatore densità, e chiedersi non più ``il metodo funziona?'' ma
``il metodo funziona abbastanza bene da sopravvivere all'hardware reale?''.
Sette documenti, mirror diretto della serie sul dimero (Documenti 5--11):
cinque stadi della pipeline di rumore, più due estensioni facoltative già
chiuse. Riusa senza modifiche il modello di rumore, l'ansatz canonico
($W$-2q.6) e i punti di lavoro stabiliti nella Parte 1 di questa stessa
topologia.
\end{abstract}

\section{Perché un sistema aperto}

Tutta la Parte 1 (Documenti 0--4) è a simulazione ideale: vettore di stato
esatto, nessun errore di gate, nessun conteggio finito di shot. È il
prerequisito corretto --- serve sapere che il metodo dà la risposta giusta
prima di chiedersi quanto rumore tollera --- ma non è la domanda che conta
per l'hardware disponibile oggi. Questa parte introduce due canali di
rumore (depolarizzante di gate, readout), calibrati su un chip reale, e
misura come si degradano gli stessi quattro strumenti già costruiti:
stato fondamentale (VQE), evoluzione temporale (Trotter), correlatori
dinamici.

\messaggio{Il rumore non è un capitolo a parte aggiunto in fondo: è la
domanda per cui tutta la Parte 1 è stata costruita con la disciplina che
ha --- verifica indipendente dove possibile (dichiarato apertamente quando
non lo è), verifiche esplicite, niente dato per buono. Qui si vede perché
quella disciplina serviva.}

\section{I sette documenti}

\begin{table}[htbp]
\centering
\footnotesize
\renewcommand{\arraystretch}{1.3}
\begin{tabular}{@{}cp{6.3cm}p{6.7cm}@{}}
\toprule
\rowcolor{grigetto}
\textbf{Doc.} & \textbf{Argomento} & \textbf{Punto di lavoro} \\
\midrule
5 & Modello di rumore: due canali, calibrazione reale & --- \\
6 & VQE sotto rumore (degradazione, non riottimizzazione) & Scenario A \\
7 & Dinamica (Trotter) sotto rumore, $N^*$ & Scenario A \emph{e} $R_0$ \\
8 & Correlatori sotto rumore, readout a shot finiti & Scenario A \\
9 & Scan sui parametri di rumore & Scenario A \\
10 & Estensione --- VQE noise-aware & Scenario A \\
11 & Estensione --- readout asimmetrico & Scenario A \\
\bottomrule
\end{tabular}
\caption{I sette documenti di questa serie. A differenza della Parte 1, dove
Scenario A e $R_0$ avevano ruoli diversi (VQE il primo, dimostrazione
Trotter il secondo), qui Scenario A è il punto principale ovunque --- $R_0$
compare solo nel Documento 7, perché è l'unico che riprende direttamente la
dimostrazione Trotter della Parte 1.}
\end{table}
\FloatBarrier

\section{Che cosa si riusa dalla Parte 1, senza modifiche}

\begin{itemize}[leftmargin=1.6em,itemsep=3pt]
\item \textbf{L'ansatz canonico}, $W$-2q.6 (Documento 2): sei parametri,
scelto perché --- a differenza di RBS --- raggiunge $\mathcal F=1$ esatto
sotto DM.
\item \textbf{Il blocco di Trotter} (Documento 3), coi suoi tre livelli
annidati (campo esatto, scambio e DM approssimati sui tre legami
condivisi).
\item \textbf{Il circuito dei correlatori} (Documento 4), quattro qubit,
ancilla al qubit 3.
\item \textbf{Il modello di rumore}, letteralmente lo stesso del dimero
(wrapper, non una riscrittura): verificato indipendente dal numero di
qubit.
\end{itemize}

\section{Una differenza di metodo rispetto alla prima versione di questa
pipeline}

Il readout, in questa serie, è sempre \textbf{genuinamente campionato} ---
un `ReadoutError` vero agganciato al simulatore, misura a shot finiti, mai
una formula analitica di correzione applicata a posteriori. È la stessa
scelta già fatta, in una revisione successiva, per la serie gemella sul
dimero: più vicina a cosa succede davvero su un dispositivo reale, dove non
si accede mai a un $\langle Z\rangle$ ideale da correggere, solo a un bit
misurato.

\messaggio{Dove la pipeline originaria di questa topologia (Parte 2,
completata prima di questa serie narrativa) usava una formula analitica per
il readout, qui è stata rifatta a shot finiti --- non per sospetto di un
errore, ma per coerenza con lo standard ora in uso, e per verificare
direttamente che i numeri non cambino. Non cambiano, entro l'errore
statistico atteso: lo si vedrà nei Documenti 8--11.}

\end{document}
